\section{Introduction}
\label{sec:intro}

Time series (TS) classification (TSC) is challenging due to temporal
dependencies, non-stationarity, or scarce labelled data, often
limited by acquisition costs and privacy. Data augmentation (DA)
mitigates these issues by enriching the training set with
synthetic samples.
In practice, however, no DA method is reliable across the
heterogeneous datasets that TSC covers
\citep{iwana2021empirical, iglesias2023data}, an inconsistency
whose practical cost we detail in
Section~\ref{sec:setup:ts-da}. This motivates the search for an
augmentation method that is safe to apply by default.

Existing in-distribution generative DA methods for TSC learn to model the training distribution
faithfully but, by design, rarely produce samples that extend
beyond its empirical support, leaving the decision boundary
under-supplied near its hardest regions, precisely where
augmentation could help the most.
\begin{figure}[t]
    \centering
    \includesvg[inkscapelatex=false,width=\textwidth]{images/ASCENSION_SVG.svg}
\caption{Comparison of latent space dynamics between
state-of-the-art generative DA methods and ASCENSION.
State-of-the-art methods generate samples confined to the
learned distribution, limiting diversity (left). ASCENSION
(right) introduces a novel $\alpha$-scaling mechanism,
controlled by a parameter~$\alpha$, which progressively expands
class regions, while a contrastive loss maintains intra-class
compactness and inter-class separation.}
    \label{fig:SoTaVSASCENSION}
\end{figure}

We posit that progressively exploring under-represented regions
of the latent space by expanding the initial data distribution
can enhance DA effectiveness, including under train/test
discrepancy. To this end, we propose ASCENSION, a DA framework
that leverages a variational autoencoder (VAE) architecture and
introduces a tunable $\alpha$-scaling mechanism to gradually
explore these under-represented regions while maintaining
intra-class compactness and inter-class separation.
Figure~\ref{fig:SoTaVSASCENSION} illustrates the core concept.
This strategy yields more diverse synthetic data and improved
mean classification performance, as shown in
Section~\ref{sec:results:main}.

The key contributions of this paper are:

\begin{itemize}
    \item \textbf{Novel VAE-based DA method.} A progressive,
    controllable boundary-expansion strategy via an
    $\alpha$-scaling mechanism, with ablations mapping the
    dataset conditions where it is most beneficial.

    \item \textbf{Comprehensive benchmark \& low-risk gains.}
    ASCENSION is evaluated across 102 UCR datasets against one
    traditional (FAA) and eight generative baselines
    (Diffusion-TS, ImagenTime, KoVAE, LatentAugment, TTS-GAN,
    Time-DDPM, VaDE, MODALS). It is the only method with a
    positive mean gain on each of the three classifiers, it
    degrades the fewest datasets, and it is never significantly
    worse than any competitor, making it a low-risk DA method.

    \item \textbf{Data-driven insights.} We investigate how
    ASCENSION's components and TS features influence
    DA effectiveness, showing that $\alpha$-scaling is the
    single most impactful component.
\end{itemize}

